% GENERATED FILE. Edit canonical lessons, then run npm run generate:books.
% canonical-source-hash: fnv1a64:d7760404cbf6ee47
% canonical-lessons: TA-C371-camattuvam, TA-C371-pakupatu, TA-C371-cirupanmai, TA-C371-perumpanmai, TA-C371-otukku

\chapter{Equal Rights}
\label{ch:ta-equal-rights}
\hlchaptermodality{371}

\begin{chapteropening}
\textbf{By the end of this chapter:} \emph{I can talk about equality, discrimination and minorities.}

Complete the last lesson of chapter 371: I can talk about equality, discrimination and minorities.
\end{chapteropening}

\section[\texorpdfstring{\ta{சமத்துவம்} (camattuvam)}{சமத்துவம் (camattuvam)}]{\ta{சமத்துவம்} (camattuvam) — equality}

\label{lesson:TA-C371-camattuvam}



\begin{quote}
Before the new one: say the Tamil for unemployment, then the Tamil for to exploit.
\end{quote}

\subsection*{You'll want to know: \ta{சமத்துவம்}}
\textbf{\ta{சமத்துவம்}} — \emph{camattuvam} — \textquotedblleft{}equality\textquotedblright{}.

\begin{quote}
\ta{ஆண்களுக்கும்} \ta{பெண்களுக்கும்} \ta{இடையே} \ta{சமத்துவம்} \ta{வேண்டும்}.
\end{quote}

\begin{quote}
\emph{āṇkaḷukkum peṇkaḷukkum iṭaiyē camattuvam vēṇṭum.}
\end{quote}

\begin{quote}
There must be equality between men and women.
\end{quote}

As in \ta{முக்கியத்துவம்}, the abstract ending -\ta{த்துவம்} joins \ta{சமம்}, equal.

\begin{cousinweb}
\textbf{\ta{சமத்துவம்}} (\emph{camattuvam}) is borrowed from Sanskrit \emph{samatva}, \textquotedblleft{}equality\textquotedblright{} (Monier-Williams), from \emph{sama}, \textquotedblleft{}equal\textquotedblright{}, and the abstract suffix \emph{-tva}.
\end{cousinweb}

\begin{grammarlens}[title={What you've built}]
One more word for giving reasons and opinions.
\end{grammarlens}

\subsection*{Your turn}
\begin{itemize}
\raggedright
  \item \emph{Say it:} \emph{camattuvam}
  \item \emph{Say it:} \emph{camattuvam}, once more
  \item \emph{Say it:} \emph{āṇkaḷukkum peṇkaḷukkum iṭaiyē camattuvam vēṇṭum.} — There must be equality between men and women.
  \item \emph{Recall:} say the Tamil for unemployment, then the Tamil for to exploit, then say \emph{camattuvam} again
\end{itemize}

\begin{tcolorbox}[breakable,colback=teal!4,colframe=teal!35!black,title={Before you move on}]
What does \ta{சமத்துவம்} mean? (\textquotedblleft{}Equality\textquotedblright{}.) Say it once more.
\end{tcolorbox}
\section[\texorpdfstring{\ta{பாகுபாடு} (pākupāṭu)}{பாகுபாடு (pākupāṭu)}]{\ta{பாகுபாடு} (pākupāṭu) — discrimination}

\label{lesson:TA-C371-pakupatu}



\begin{quote}
Before the new one: say the Tamil for to exploit, then the Tamil for equality.
\end{quote}

\subsection*{You'll want to know: \ta{பாகுபாடு}}
\textbf{\ta{பாகுபாடு}} — \emph{pākupāṭu} — \textquotedblleft{}discrimination\textquotedblright{}.

\begin{quote}
\ta{வேலையில்} \ta{யாரிடமும்} \ta{பாகுபாடு} \ta{காட்டக்} \ta{கூடாது}.
\end{quote}

\begin{quote}
\emph{vēlaiyil yāriṭamum pākupāṭu kāṭṭak kūṭātu.}
\end{quote}

\begin{quote}
Nobody should be discriminated against at work.
\end{quote}

\ta{பாகுபாடு} \ta{காட்டு}, to discriminate.

\begin{cousinweb}
\textbf{\ta{பாகுபாடு}} (\emph{pākupāṭu}) is built from \emph{pāku}, \textquotedblleft{}a portion, a division\textquotedblright{}, usually traced to Sanskrit \emph{bhāga}, \textquotedblleft{}portion\textquotedblright{} (CDIAL 9430), and \emph{pāṭu}, the noun of \emph{paṭu}: a dividing of people into parts.
\end{cousinweb}

\begin{grammarlens}[title={What you've built}]
One more word for giving reasons and opinions.
\end{grammarlens}

\subsection*{Your turn}
\begin{itemize}
\raggedright
  \item \emph{Say it:} \emph{pākupāṭu}
  \item \emph{Say it:} \emph{pākupāṭu}, once more
  \item \emph{Say it:} \emph{vēlaiyil yāriṭamum pākupāṭu kāṭṭak kūṭātu.} — Nobody should be discriminated against at work.
  \item \emph{Recall:} say the Tamil for to exploit, then the Tamil for equality, then say \emph{pākupāṭu} again
\end{itemize}

\begin{tcolorbox}[breakable,colback=teal!4,colframe=teal!35!black,title={Before you move on}]
What does \ta{பாகுபாடு} mean? (\textquotedblleft{}Discrimination\textquotedblright{}.) Say it once more.
\end{tcolorbox}
\section[\texorpdfstring{\ta{சிறுபான்மை} (ciṟupāṉmai)}{சிறுபான்மை (ciṟupāṉmai)}]{\ta{சிறுபான்மை} (ciṟupāṉmai) — a minority}

\label{lesson:TA-C371-cirupanmai}



\begin{quote}
Before the new one: say the Tamil for equality, then the Tamil for discrimination.
\end{quote}

\subsection*{You'll want to know: \ta{சிறுபான்மை}}
\textbf{\ta{சிறுபான்மை}} — \emph{ciṟupāṉmai} — \textquotedblleft{}a minority\textquotedblright{}.

\begin{quote}
\ta{இந்த} \ta{ஊரில்} \ta{அவர்கள்} \ta{சிறுபான்மையாக} \ta{இருக்கிறார்கள்}.
\end{quote}

\begin{quote}
\emph{inta ūril avarkaḷ ciṟupāṉmaiyāka irukkiṟārkaḷ.}
\end{quote}

\begin{quote}
In this town they are a minority.
\end{quote}

\ta{சிறு}, small + \ta{பான்மை}, portion.

\begin{cousinweb}
\textbf{\ta{சிறுபான்மை}} (\emph{ciṟupāṉmai}) is native: \emph{ciṟu}, \textquotedblleft{}small\textquotedblright{} (DEDR 1594), and \emph{pāṉmai}, \textquotedblleft{}portion, share\textquotedblright{} (DEDR 4097): the small portion.
\end{cousinweb}

\begin{grammarlens}[title={What you've built}]
One more word for giving reasons and opinions.
\end{grammarlens}

\subsection*{Your turn}
\begin{itemize}
\raggedright
  \item \emph{Say it:} \emph{ciṟupāṉmai}
  \item \emph{Say it:} \emph{ciṟupāṉmai}, once more
  \item \emph{Say it:} \emph{inta ūril avarkaḷ ciṟupāṉmaiyāka irukkiṟārkaḷ.} — In this town they are a minority.
  \item \emph{Recall:} say the Tamil for equality, then the Tamil for discrimination, then say \emph{ciṟupāṉmai} again
\end{itemize}

\begin{tcolorbox}[breakable,colback=teal!4,colframe=teal!35!black,title={Before you move on}]
What does \ta{சிறுபான்மை} mean? (\textquotedblleft{}A minority\textquotedblright{}.) Say it once more.
\end{tcolorbox}
\section[\texorpdfstring{\ta{பெரும்பான்மை} (perumpāṉmai)}{பெரும்பான்மை (perumpāṉmai)}]{\ta{பெரும்பான்மை} (perumpāṉmai) — a majority}

\label{lesson:TA-C371-perumpanmai}



\begin{quote}
Before the new one: say the Tamil for discrimination, then the Tamil for a minority.
\end{quote}

\subsection*{You'll want to know: \ta{பெரும்பான்மை}}
\textbf{\ta{பெரும்பான்மை}} — \emph{perumpāṉmai} — \textquotedblleft{}a majority\textquotedblright{}.

\begin{quote}
\ta{பெரும்பான்மை} \ta{மாணவர்கள்} \ta{இந்த} \ta{யோசனையை} \ta{ஏற்றுக்கொண்டார்கள்}.
\end{quote}

\begin{quote}
\emph{perumpāṉmai māṇavarkaḷ inta yōcaṉaiyai ēṟṟukkoṇṭārkaḷ.}
\end{quote}

\begin{quote}
The majority of the students accepted this idea.
\end{quote}

\ta{பெரும்பான்மையான}, most: \ta{பெரும்பான்மையான} \ta{மக்கள்}, most people.

\begin{cousinweb}
\textbf{\ta{பெரும்பான்மை}} (\emph{perumpāṉmai}) is native: \emph{perum}, \textquotedblleft{}great\textquotedblright{}, from \emph{peru} (DEDR 4411), and \emph{pāṉmai}, \textquotedblleft{}portion, share\textquotedblright{} (DEDR 4097): the greater portion.
\end{cousinweb}

\begin{grammarlens}[title={What you've built}]
One more word for giving reasons and opinions.
\end{grammarlens}

\subsection*{Your turn}
\begin{itemize}
\raggedright
  \item \emph{Say it:} \emph{perumpāṉmai}
  \item \emph{Say it:} \emph{perumpāṉmai}, once more
  \item \emph{Say it:} \emph{perumpāṉmai māṇavarkaḷ inta yōcaṉaiyai ēṟṟukkoṇṭārkaḷ.} — The majority of the students accepted this idea.
  \item \emph{Recall:} say the Tamil for discrimination, then the Tamil for a minority, then say \emph{perumpāṉmai} again
\end{itemize}

\begin{tcolorbox}[breakable,colback=teal!4,colframe=teal!35!black,title={Before you move on}]
What does \ta{பெரும்பான்மை} mean? (\textquotedblleft{}A majority\textquotedblright{}.) Say it once more.
\end{tcolorbox}
\section[\texorpdfstring{\ta{ஒடுக்கு} (oṭukku)}{ஒடுக்கு (oṭukku)}]{\ta{ஒடுக்கு} (oṭukku) — to oppress; to suppress}

\label{lesson:TA-C371-otukku}



\begin{quote}
Before the new one: say the Tamil for a minority, then the Tamil for a majority.
\end{quote}

\subsection*{You'll want to know: \ta{ஒடுக்கு}}
\textbf{\ta{ஒடுக்கு}} — \emph{oṭukku} — \textquotedblleft{}to oppress; to suppress\textquotedblright{}.

\begin{quote}
\ta{அந்த} \ta{அரசன்} \ta{தன்} \ta{மக்களை} \ta{ஒடுக்கினான்}.
\end{quote}

\begin{quote}
\emph{anta aracaṉ taṉ makkaḷai oṭukkiṉāṉ.}
\end{quote}

\begin{quote}
That king oppressed his people.
\end{quote}

Past (I …): \ta{ஒடுக்கினேன்}, \emph{oṭukkiṉēṉ}.

Who is oppressed takes -\ta{ஐ}. The noun is \ta{ஒடுக்குமுறை}, oppression.

\begin{cousinweb}
\textbf{\ta{ஒடுக்கு}} (\emph{oṭukku}) is native Dravidian: DEDR 954 lists Tamil \emph{oṭuṅku}, \textquotedblleft{}to be restrained, shrink\textquotedblright{}, and its causative \emph{oṭukku}, \textquotedblleft{}to restrain, repress\textquotedblright{}.
\end{cousinweb}

\begin{grammarlens}[title={What you've built}]
One more word for giving reasons and opinions.
\end{grammarlens}

\subsection*{Your turn}
\begin{itemize}
\raggedright
  \item \emph{Say it:} \emph{oṭukku}
  \item \emph{Say it:} \emph{oṭukku}, once more
  \item \emph{Say it:} \emph{anta aracaṉ taṉ makkaḷai oṭukkiṉāṉ.} — That king oppressed his people.
  \item \emph{Recall:} say the Tamil for a minority, then the Tamil for a majority, then say \emph{oṭukku} again
\end{itemize}

\begin{tcolorbox}[breakable,colback=teal!4,colframe=teal!35!black,title={Before you move on}]
What does \ta{ஒடுக்கு} mean? (\textquotedblleft{}To oppress; to suppress\textquotedblright{}.) Say it once more.
\end{tcolorbox}
